	\section{Statistical comparison}
	\subsection{Energy-distribution overlap criterion}

	Beyond purely statistical metrics, the \emph{physical content} of
	competing models can be compared through their implied site-energy
	distributions. Each model that admits an energy distribution
	$f(E)$ is characterised by a mean energy $E_c$ and a standard
	deviation $\sigma_E$. The central idea is: \textbf{if the energy
	distributions of two models overlap to the point that they are
	statistically indistinguishable, the simpler model should be
	preferred.}

	\paragraph{Summary of energy parameters by model.}

	\begin{itemize}
		\item \textbf{Langmuir}: Dirac delta at
		$E_0 = RT\ln(K_L C^0)$, with $\sigma_E = 0$.
		\item \textbf{Double Langmuir}: two Dirac deltas at
		$E_1 = RT\ln(K_1 C^0)$ and $E_2 = RT\ln(K_2 C^0)$,
		separated by
		\begin{equation}
			\Delta E_{12} = |E_1 - E_2|
			= RT\left|\ln\frac{K_1}{K_2}\right|.
			\label{eq:DeltaE12}
		\end{equation}
		\item \textbf{Sips}: logistic distribution centred at
		$E_c = RT\ln(K_S C^{\ast})$ with standard deviation
		\begin{equation}
			\sigma_E^{\mathrm{(Sips)}}
			= \frac{\pi}{\sqrt{3}}\,nRT,
			\label{eq:sigma_sips_stat}
		\end{equation}
		where $n$ is the Sips heterogeneity parameter (note: $1/n = m$
		is the exponent in Eq.~\eqref{eq:Sips}).
		\item \textbf{Toth}: asymmetric distribution with mean energy
		$E_c$ and a skewness-dependent width; the effective spread is
		broader toward low energies.
		\item \textbf{D-R/D-A}: Gaussian-like distribution with
		$\sigma_E \propto E_a / \sqrt{2}$ (D-R) or dependent on
		$n_{DA}$ (D-A).
	\end{itemize}

	\paragraph{Criterion 1: Langmuir vs.\ Sips.}

	The Sips model reduces to Langmuir when $n \to 1$
	(or $m \to 1$), i.e.\ when the energy distribution collapses to
	a delta. The question is: given the fitted value of $n$, is the
	distribution narrow enough that the heterogeneity is negligible?

	A natural criterion compares the Sips distribution width with the
	thermal energy scale:
	\begin{equation}
		\boxed{
			\text{Accept Langmuir if}\quad
			\sigma_E^{\mathrm{(Sips)}}
			= \frac{\pi}{\sqrt{3}}\,nRT
			\leq \alpha\, RT,
		}
		\label{eq:criterion_L_vs_S}
	\end{equation}
	where $\alpha$ is a tolerance (typically $\alpha \approx 1$--$2$).
	This corresponds to requiring $n \leq \alpha\sqrt{3}/\pi$, i.e.\
	$n \lesssim 0.55$--$1.10$. In practice, the F-test
	(Eq.~\eqref{eq:Ftest}) already provides a frequentist check;
	Eq.~\eqref{eq:criterion_L_vs_S} adds a \emph{physical}
	interpretation: the additional parameter in Sips is justified only
	if the implied energy spread exceeds the thermal fluctuation scale.

	\paragraph{Criterion 2: Langmuir vs.\ Double Langmuir.}

	The Double Langmuir model introduces two distinct site types.
	The two sites are physically distinguishable only if their
	energy separation exceeds the thermal broadening:
	\begin{equation}
		\boxed{
			\text{Accept Double Langmuir if}\quad
			\Delta E_{12}
			= RT\left|\ln\frac{K_1}{K_2}\right|
			> \beta\, RT,
		}
		\label{eq:criterion_L_vs_DL}
	\end{equation}
	where $\beta$ is a threshold (typically $\beta \approx 2$--$3$,
	meaning the site energies differ by at least $2$--$3\,RT$).
	If $\Delta E_{12} \leq \beta\,RT$, the two ``sites'' are
	thermally indistinguishable and a single Langmuir model is
	more parsimonious.

	Equivalently, in terms of the ratio of Langmuir constants:
	\begin{equation}
		\frac{K_1}{K_2} > \mathrm{e}^{\beta}
		\quad\Longleftrightarrow\quad
		\Delta E_{12} > \beta\, RT.
	\end{equation}
	For $\beta = 2$: $K_1/K_2 > 7.4$. For $\beta = 3$:
	$K_1/K_2 > 20$.

	\paragraph{Criterion 3: Double Langmuir vs.\ Sips.}

	When both the Double Langmuir and the Sips model fit the data
	well, one can compare the two-delta representation (DL) with the
	continuous logistic distribution (Sips). The Sips model is
	preferred if the continuous distribution adequately covers the
	two peaks:
	\begin{equation}
		\boxed{
			\text{Prefer Sips over DL if}\quad
			\Delta E_{12} < 2\,\sigma_E^{\mathrm{(Sips)}}.
		}
		\label{eq:criterion_DL_vs_S}
	\end{equation}
	Conversely, if $\Delta E_{12} \gg 2\sigma_E^{\mathrm{(Sips)}}$,
	the two-site model captures bimodal heterogeneity that the
	symmetric Sips distribution cannot resolve.

	\paragraph{Criterion 4: general overlap test.}

	For any two models $A$ and $B$ that each imply an energy
	distribution $f_A(E)$ and $f_B(E)$ with means $E_c^{A}$,
	$E_c^{B}$ and standard deviations $\sigma_A$, $\sigma_B$,
	the overlap coefficient
	\begin{equation}
		\mathrm{OVL}
		= \int_{-\infty}^{\infty}
		\min\!\bigl[\hat{f}_A(E),\;\hat{f}_B(E)\bigr]\,\mathrm{d}E,
		\label{eq:OVL}
	\end{equation}
	where $\hat{f}$ denotes the normalised densities, quantifies the
	degree of agreement. $\mathrm{OVL} = 1$ means the distributions
	are identical; $\mathrm{OVL} \approx 0$ means they are well
	separated.

	If both distributions are approximately Gaussian (a good
	approximation for Sips, D-R, and Toth near the peak), the overlap
	has the closed-form estimate
	\begin{equation}
		\mathrm{OVL}
		\approx 2\,\Phi\!\left(
		-\frac{|E_c^A - E_c^B|}{2\,\bar{\sigma}}
		\right),
		\qquad
		\bar{\sigma}
		= \sqrt{\frac{\sigma_A^2 + \sigma_B^2}{2}},
		\label{eq:OVL_gauss}
	\end{equation}
	where $\Phi$ is the standard normal CDF. The rule of thumb:
	\begin{itemize}
		\item $\mathrm{OVL} > 0.8$: models are practically
		indistinguishable --- prefer the simpler one.
		\item $0.3 < \mathrm{OVL} < 0.8$: partial overlap; use
		AIC/BIC to decide.
		\item $\mathrm{OVL} < 0.3$: distributions are well separated;
		the models describe qualitatively different physics.
	\end{itemize}

	\subsection{Classical statistical comparison}

	\paragraph{Two models.}
	For two nested models (e.g.\ Langmuir vs.\ Double Langmuir),
	the F-test (Eq.~\eqref{eq:Ftest}) is the standard approach.
	For two non-nested models with the same $p$ (e.g.\ Langmuir vs.\
	Freundlich), direct comparison of AIC or BIC is appropriate.
	The Akaike weight of model $i$ among a set of $M$ candidates is
	\begin{equation}
		w_i = \frac{\exp(-\Delta\mathrm{AIC}_i / 2)}
		{\displaystyle\sum_{j=1}^{M}
			\exp(-\Delta\mathrm{AIC}_j / 2)},
		\label{eq:Akaike_weight}
	\end{equation}
	and gives the probability that model $i$ is the best model given
	the data.

	\paragraph{Several models.}
	When comparing $M > 2$ models simultaneously:
	\begin{enumerate}
		\item Compute AIC$_c$ for each model.
		\item Rank by $\Delta\mathrm{AIC}$ and compute Akaike weights
		(Eq.~\eqref{eq:Akaike_weight}).
		\item Discard models with $\Delta\mathrm{AIC} > 10$.
		\item If multiple models remain ($\Delta\mathrm{AIC} < 2$),
		report all as plausible and use model-averaged parameters:
		\begin{equation}
			\hat{\theta}_{\mathrm{avg}}
			= \sum_{i=1}^{M} w_i\,\hat{\theta}_i.
			\label{eq:model_averaging}
		\end{equation}
	\end{enumerate}
	Complementary visualisation: error-metric heatmaps
	($\log_{10}$-transformed) across substrates and models allow
	rapid identification of the best model for each material.

	\subsubsection{Combined decision protocol}

	We recommend the following workflow for model selection:
	\begin{enumerate}
		\item Fit all candidate models by non-linear regression.
		\item Compute residual metrics (SSE, ARED, HYBRID) and
		information criteria (AIC$_c$, BIC).
		\item Apply the F-test for nested pairs.
		\item Compute the energy-distribution parameters ($E_c$,
		$\sigma_E$, or $\Delta E_{12}$) from the fitted constants.
		\item Apply the overlap criteria
		(Eqs.~\eqref{eq:criterion_L_vs_S}--\eqref{eq:OVL_gauss}):
		if a more complex model is statistically preferred (step~3)
		but the energy distributions are indistinguishable
		($\mathrm{OVL} > 0.8$), retain the simpler model and flag
		the result as ``statistically significant but physically
		negligible heterogeneity''.
		\item Report both the statistical ranking \emph{and} the
		energy-distribution comparison to give the reader a complete
		picture.
	\end{enumerate}

